\documentclass[11pt,a4paper]{article}
\usepackage[margin=25mm,headheight=15pt]{geometry}
\usepackage[T1]{fontenc}
\usepackage[utf8]{inputenc}
\usepackage{lmodern,amsmath,amssymb,amsthm,mathtools,bm}
\usepackage{xcolor,graphicx,booktabs,tabularx,array,longtable}
\usepackage{microtype,enumitem,fancyhdr}
\usepackage{xurl}
\usepackage{hyperref}
\definecolor{accent}{HTML}{163B55}
\hypersetup{colorlinks=true,linkcolor=accent,citecolor=accent,urlcolor=accent,
 pdftitle={A strict numerical-index bound for real Lebesgue spaces},pdfauthor={Matthew J. Colbrook},
 pdfsubject={AIM 935: mathematical result and proof}}
\urlstyle{same}
\setlength{\parindent}{0pt}
\setlength{\parskip}{0.35em}
\setlength{\emergencystretch}{3em}
\setlist{itemsep=0.25em,topsep=0.3em}
\newtheorem{theorem}{Theorem}[section]
\newtheorem{lemma}[theorem]{Lemma}
\newtheorem{proposition}[theorem]{Proposition}
\newtheorem{corollary}[theorem]{Corollary}
\theoremstyle{remark}
\newtheorem{remark}[theorem]{Remark}
\newtheorem{definition}[theorem]{Definition}
\DeclareMathOperator{\Tr}{Tr}
\DeclareMathOperator{\supp}{supp}
\DeclareMathOperator{\Lip}{Lip}
\DeclareMathOperator{\Rea}{Re}
\newcommand{\R}{\mathbb R}
\newcommand{\C}{\mathbb C}
\newcommand{\E}{\mathbb E}
\newcommand{\N}{\mathbb N}
\newcommand{\ip}[2]{\langle #1,#2\rangle}
\pagestyle{fancy}
\fancyhf{}
\fancyhead[L]{\small\sffamily AIM 935}
\fancyhead[R]{\small\sffamily Matthew J. Colbrook}
\fancyfoot[C]{\thepage}
\renewcommand{\headrulewidth}{0.3pt}
\allowdisplaybreaks[2]
\title{\textbf{A strict numerical-index bound for real Lebesgue spaces}}
\author{Matthew J. Colbrook\\[3pt]
{\small Department of Applied Mathematics and Theoretical Physics}\\
{\small University of Cambridge, United Kingdom}\\
{\small\href{mailto:mjc249@maths.cam.ac.uk}{\texttt{mjc249@maths.cam.ac.uk}}}}
\date{3 November 2026}
\begin{document}
\maketitle
\begin{abstract}
We show that the numerical index of a real Lebesgue space is strictly below the proposed two-dimensional value at every non-Hilbert exponent between one and infinity. A three-dimensional analytic perturbation proves the general strict inequality. Independent exact algebraic and rational examples give explicit certificates at exponent four and its conjugate.
\end{abstract}
\noindent\textbf{Repository target.}
This manuscript concerns AIM problem 935 at revision
\texttt{5d7e26e8444e}. The precise scope is stated below.
\medskip

\section{Main result and proof}
\label{sec:935}

\paragraph{Answer.} The proposed identity is false for every
$1<p<\infty$ with $p\ne2$. A three-dimensional perturbation gives the
strict inequality throughout this range. An independent rational
certificate at $p=4$ supplies a concrete matrix and an explicit numerical
gap; its transpose also treats $p=4/3$.

\subsection{Strict inequality for every non-Hilbert exponent}
\label{sec:935-general}

The counterexample admits an analytic extension to every exponent in the
question. The explicit rational certificate below remains useful because
it supplies a concrete matrix and a numerical gap at $p=4$.

\begin{theorem}\label{thm:935-allp}
For every $1<p<\infty$ with $p\ne2$,
\[
 n\bigl(L_p([0,1];\mathbb R)\bigr)
 \le n\bigl(\ell_p^3(\mathbb R)\bigr)<M_p.
\]
In particular, the proposed equality fails for every exponent to which
Problem~935 applies.
\end{theorem}

We prove the theorem for $p>2$ by perturbing an operator that attains the
candidate ratio. The perturbation changes its numerical radius by at most
$O(\varepsilon^2)$, while increasing a lower bound for its norm by a
positive multiple of $|\varepsilon|^{p/(p-1)}$. The latter exponent is
strictly smaller than two.

\begin{lemma}\label{lem:935-scalar}
For $p>2$, the function
\[
 f_p(t)=\frac{t-t^{p-1}}{1+t^p},\qquad 0\le t\le1,
\]
has a unique maximum at some $t=t_p\in(0,1)$, and $f_p''(t_p)<0$.
Set $j_p(s)=|s|^{p-2}s$ for real $s$, with $j_p(0)=0$. The function
\[
 h_p(s)=\frac{j_p(s)-s}{1+|s|^p}
\]
has maximum $M_p$ precisely at $s=t_p^{-1}$ and $s=-t_p$, and minimum
$-M_p$ precisely at $s=t_p$ and $s=-t_p^{-1}$. Each of these extrema is
nondegenerate.
\end{lemma}
\begin{proof}
Write $t=e^{-x}$, $a=(p-2)/2$, and $b=p/2$. For $x>0$,
\[
 f_p(e^{-x})=\frac{\sinh(ax)}{\cosh(bx)}.
\]
Its logarithmic derivative is
\[
 a\coth(ax)-b\tanh(bx),
\]
whose derivative is
\[
 -a^2\operatorname{csch}^2(ax)-b^2\operatorname{sech}^2(bx)<0.
\]
The logarithmic derivative tends to $+\infty$ as $x\downarrow0$ and to
$a-b=-1$ as $x\to\infty$. It therefore has exactly one zero, where the
second derivative of the positive function is strictly negative. The
change of variable transfers nondegeneracy to $t_p$. Both endpoint values
of $f_p$ are zero, so this is its unique global maximum.
The assertions about $h_p$ follow from its oddness and the relations
$h_p(s)=-f_p(s)$ for $0<s<1$ and $h_p(s)=f_p(1/s)$ for $s>1$.
\end{proof}

\begin{proof}[Proof of Theorem~\ref{thm:935-allp} for $p>2$]
Fix $p>2$, write $q=p/(p-1)$, and put
\[
 c=2^{-1/p},\qquad u=(c,c),\qquad w=(c,-c).
\]
Thus $\|u\|_p=\|w\|_p=1$ and
$\langle j_p(u),w\rangle=0$. Regard $\ell_p^3$ as
$\ell_p^2\oplus_p\mathbb R$, with coordinates $(\xi,y)$, where
$\xi=(x,z)$. Define
\[
 T_0(\xi,y)=\bigl(yu,-\langle j_p(u),\xi\rangle\bigr).
\]
Since $\|j_p(u)\|_q=1$,
\[
 \|T_0(\xi,y)\|_p^p
 =|y|^p+|\langle j_p(u),\xi\rangle|^p
 \le |y|^p+\|\xi\|_p^p.
\]
Equality occurs, for example, when $\xi=0$, so $\|T_0\|_p=1$.

For any real matrix $T$, smoothness of the $\ell_p$ norm gives
\[
 v_p(T)=\sup_{v\ne0}
 \frac{|\langle j_p(v),Tv\rangle|}{\|v\|_p^p}.
 \tag{935.G1}\label{eq:935-general-radius}
\]
Set $g_p(s,t)=j_p(s)t-j_p(t)s$. The scalar inequality defining $M_p$
implies
\[
 |g_p(s,t)|\le M_p(|s|^p+|t|^p)
 \qquad(s,t\in\mathbb R).
\]
The numerical form of $T_0$ is
\[
 F_0(x,z,y)=g_p(x,cy)+g_p(z,cy).
\]
Because $2c^p=1$, summing the preceding scalar inequalities yields
$|F_0(x,z,y)|\le M_p(|x|^p+|z|^p+|y|^p)$. Equality is obtained by choosing
both scalar pairs to attain the same signed maximum. Consequently
\[
 v_p(T_0)=M_p.
 \tag{935.G2}\label{eq:935-star-radius}
\]

We next describe all maximizing vectors and their local behavior.
A maximizing vector has $y\ne0$: otherwise $F_0=0$ while $M_p>0$.
After a common nonzero real rescaling, use the chart $cy=1$ and set
\[
 S=|x|^p+|z|^p+2.
\]
In this chart the normalized numerical form is
\[
 \phi_0(x,z)
 =\frac{(1+|x|^p)h_p(x)+(1+|z|^p)h_p(z)}{S}.
 \tag{935.G3}\label{eq:935-star-chart}
\]
Every weight is positive. Thus $\phi_0=M_p$ exactly when
\[
 x,z\in\{t_p^{-1},-t_p\},
\]
and $\phi_0=-M_p$ exactly when
\[
 x,z\in\{t_p,-t_p^{-1}\}.
\]
These choices form a finite collection of points. Lemma~\ref{lem:935-scalar} also proves
that each maximum of $|\phi_0|$ is nondegenerate. For example, at a positive
maximum,
\[
 M_p-\phi_0(x,z)
 =\frac{(1+|x|^p)(M_p-h_p(x))
       +(1+|z|^p)(M_p-h_p(z))}{S}.
\]
Each one-variable deficit has a zero of order exactly two, with positive
second derivative. The Hessian of the displayed deficit at a zero is
therefore diagonal and positive definite. Negative maxima are handled
by replacing $h_p$ with $-h_p$.

Let $t=t_p$ and define
\[
 \kappa=\frac{1-t^{2p}}{t^2-t^{2p-2}},\qquad
 R=\begin{pmatrix}1&\kappa\\-\kappa&-1\end{pmatrix}.
\]
The denominator is positive, and
\[
 \kappa-1
 =\frac{(1-t^2)(1+t^{2p-2})}{t^2-t^{2p-2}}>0.
 \tag{935.G4}\label{eq:935-kappa}
\]
For real $\varepsilon$ set
\[
 T_\varepsilon(\xi,y)
 =\bigl(yu+\varepsilon R\xi,-\langle j_p(u),\xi\rangle\bigr).
 \tag{935.G5}\label{eq:935-family}
\]
Its numerical form is $F_0+\varepsilon G$, where
\[
 G(x,z)=|x|^p-|z|^p
       +\kappa\bigl(j_p(x)z-j_p(z)x\bigr).
 \tag{935.G6}\label{eq:935-perturbation}
\]
This perturbation vanishes at every maximizing vector of $T_0$.
Indeed, their leaf-coordinate ratios are
$x/z\in\{1,-t^2,-t^{-2}\}$. The first ratio gives zero directly. For the
second ratio,
\[
 G(-t^2,1)=t^{2p}-1+\kappa(t^2-t^{2p-2})=0.
\]
The third follows from antisymmetry under exchanging $x,z$ and homogeneity.

It follows that there are finite constants $C\ge0$ and $\varepsilon_0>0$
(depending on $p$) such that
\[
 v_p(T_\varepsilon)\le M_p+C\varepsilon^2
 \qquad(|\varepsilon|\le\varepsilon_0).
 \tag{935.G7}\label{eq:935-quadratic}
\]
Here are the compactness details. On the unit sphere, the maximizing set
of $|F_0|$ is finite. In a chart about each such point, nondegeneracy gives
$|\phi_0|\le M_p-c_0\|h\|_2^2$ for some $c_0>0$. The normalized
perturbation is smooth in that chart, because all coordinates of the
maximizer are nonzero, and vanishes at its center. It is therefore bounded
in absolute value by $L_0\|h\|_2$. Completing the square gives
\[
 |\phi_0+\varepsilon G/S|
 \le M_p-c_0\|h\|_2^2+|\varepsilon|L_0\|h\|_2
 \le M_p+\frac{L_0^2}{4c_0}\varepsilon^2.
\]
Outside the finitely many neighborhoods, compactness gives a fixed strict
gap below $M_p$, while the normalized perturbation is bounded. For
sufficiently small $|\varepsilon|$ this region contributes at most $M_p$.
Taking the largest of the finitely many constants proves
\eqref{eq:935-quadratic}.

The norm has a larger leading-order increase. Since
$Rw=(1-\kappa)u$ and $\langle j_p(u),w\rangle=0$,
\[
 T_\varepsilon(sw,y)
 =\bigl((y-(\kappa-1)\varepsilon s)u,0\bigr).
\]
The input norm is $(|s|^p+|y|^p)^{1/p}$. Taking the norm of this scalar
functional on $\ell_p^2$ gives
\[
 \|T_\varepsilon\|_p
 \ge\bigl(1+(\kappa-1)^q|\varepsilon|^q\bigr)^{1/q}.
 \tag{935.G8}\label{eq:935-general-norm}
\]
As $1<q<2$, equations \eqref{eq:935-quadratic} and
\eqref{eq:935-general-norm} imply
\[
 \begin{split}
 \frac{v_p(T_\varepsilon)}{\|T_\varepsilon\|_p}
 &\le\frac{M_p+C\varepsilon^2}
 {\bigl(1+(\kappa-1)^q|\varepsilon|^q\bigr)^{1/q}}\\
 &=M_p-\frac{M_p}{q}(\kappa-1)^q|\varepsilon|^q
      +o(|\varepsilon|^q).
 \end{split}
\]
The leading correction is strictly negative. Hence every sufficiently
small nonzero $\varepsilon$ gives
$v_p(T_\varepsilon)/\|T_\varepsilon\|_p<M_p$, proving
$n(\ell_p^3(\mathbb R))<M_p$.
\end{proof}

\paragraph{Dual exponents and the Lebesgue space.}
For any operator $T$ on a finite-dimensional real normed space,
$\|T^*\|=\|T\|$ and $v(T^*)=v(T)$. Indeed, the correspondence between
norming pairs $(x,x^*)$ and $(x^*,Jx)$ preserves the scalar $x^*(Tx)$,
and finite-dimensional reflexivity makes the correspondence exhaustive.
Thus the transpose of a counterexample on $\ell_p^3$ gives one on
$\ell_q^3$, where $q=p/(p-1)$. The substitution $s=t^{p-1}$ in the maximum
shows $M_q=M_p$. This proves the strict inequality for $1<p<2$ as well.

Finally, for any real $3\times3$ matrix $T$, its pointwise action on
$L_p([0,1];\ell_p^3)$ has the same operator norm and numerical radius as
$T$. The upper bounds follow by integrating their pointwise defining
inequalities, and constant vector-valued functions give the reverse
inequalities. This vector-valued space is isometric to scalar real
$L_p([0,1])$ by sending the three coordinates to three equal subintervals
and multiplying values by $3^{1/p}$. Consequently
$n(L_p([0,1];\mathbb R))\le n(\ell_p^3(\mathbb R))$, and the proof of
Theorem~\ref{thm:935-allp} is complete.


\subsection{A short explicit certificate at \texorpdfstring{$p=4$}{p=4}}
\label{sec:935-clean}

The perturbation family also gives an explicit algebraic example with a
short direct proof, independent of the compactness argument. Put
\[
 c=2^{-1/4},\qquad \varepsilon=2^{-12},\qquad
 T=\begin{pmatrix}
 \varepsilon&4\varepsilon&c\\
 -4\varepsilon&-\varepsilon&c\\
 -c^3&-c^3&0
 \end{pmatrix}
 \quad\hbox{on real }\ell_4^3,
 \tag{935.C1}\label{eq:935-clean-matrix}
\]
with coordinates ordered as $(x,z,y)$. We prove
\[
 \frac{v_4(T)}{\|T\|_4}<M_4=\frac1{2\sqrt2}.
\]

Write $b=cy$, $S=x^4+z^4+y^4=x^4+z^4+2b^4$, and
\[
 F_0=b(x^3+z^3)-b^3(x+z),\qquad
 G=(x^2-z^2)(x^2+z^2+4xz).
\]
The numerical form of $T$ is $F_0+\varepsilon G$. Define
\[
 Q_x=x^2-\sqrt2\,bx-b^2,\qquad
 Q_z=z^2-\sqrt2\,bz-b^2.
\]
Two direct polynomial identities are
\[
 M_4S-F_0=M_4(Q_x^2+Q_z^2),
 \tag{935.C2}\label{eq:935-clean-deficit}
\]
\[
 G=(x+z)\bigl[(x+3z+\sqrt2 b)Q_x
                   -(3x+z+\sqrt2 b)Q_z\bigr].
 \tag{935.C3}\label{eq:935-clean-G}
\]

On $S=1$, the Euclidean norm of the coefficient vector multiplying
$(Q_x,Q_z)$ in \eqref{eq:935-clean-G} is at most
\[
 |x+z|\bigl(4\|(x,z)\|_2+2|b|\bigr)\le20^{3/4}.
 \tag{935.C4}\label{eq:935-clean-coefficients}
\]
For completeness, let $r=(x^4+z^4)^{1/4}$. Then
$|x+z|\le2^{3/4}$, $\|(x,z)\|_2\le2^{1/4}r$, and
$2|b|=2^{3/4}|y|$, while $r^4+|y|^4=1$. H\"older's inequality gives
\[
 4\|(x,z)\|_2+2|b|
 \le\bigl((4\,2^{1/4})^{4/3}+(2^{3/4})^{4/3}\bigr)^{3/4}
 =(8+2)^{3/4}.
\]
This proves \eqref{eq:935-clean-coefficients}.

Set $r_Q=(Q_x^2+Q_z^2)^{1/2}$. Equations
\eqref{eq:935-clean-deficit}--\eqref{eq:935-clean-coefficients} yield
\[
 \begin{split}
 F_0+\varepsilon G
 &\le M_4-M_4r_Q^2+|\varepsilon|20^{3/4}r_Q\\
 &\le M_4+\frac{20^{3/2}}{4M_4}\varepsilon^2
  =M_4+20\sqrt{10}\,\varepsilon^2.
 \end{split}
\]
Replacing $b$ by $-b$ and $\varepsilon$ by $-\varepsilon$ gives the same argument for
$-F_0-\varepsilon G$. Hence
\[
 v_4(T)\le M_4+20\sqrt{10}\,\varepsilon^2.
 \tag{935.C5}\label{eq:935-clean-v}
\]

For the norm bound, put $u=(c,c)$ and $w=(c,-c)$. On the isometric copy of
$\ell_4^2$ consisting of vectors $(sw,y)$,
\[
 T(sw,y)=\bigl((y-3\varepsilon s)u,0\bigr).
\]
The norm of this scalar functional is therefore a lower bound for the full
operator norm:
\[
 \|T\|_4\ge\bigl(1+3^{4/3}|\varepsilon|^{4/3}\bigr)^{3/4}.
 \tag{935.C6}\label{eq:935-clean-norm}
\]
Let $a=3^{4/3}2^{-16}\in(0,1)$. The elementary inequality
$(1+a)^{3/4}\ge1+3a/8$ on $[0,1]$ follows, for example, by integrating
$(3/4)(1+a)^{-1/4}\ge3/8$. Thus the upper bound in
\eqref{eq:935-clean-v} is strictly less than $M_4$ times the lower bound in
\eqref{eq:935-clean-norm}: after division by $M_4$ and multiplication by
$2^{24}$, the required comparison is exactly
\[
 80\sqrt5<96\,3^{4/3}.
\]
It follows already from
$80\sqrt5<240<288<96\,3^{4/3}$. This proves the claimed strict ratio.
Pointwise amplification transfers this explicit example to real
$L_4([0,1])$.


\subsection{An explicit rational certificate at \texorpdfstring{$p=4$}{p=4}}

\begin{theorem}\label{thm:935}
For $p\in\{4,4/3\}$,
\[
 n\bigl(L_p([0,1];\mathbb R)\bigr)
 <\frac{1667}{4716}
 <\frac1{2\sqrt2}=M_p.
\]
The rational upper bound is approximately $0.35347752$, whereas
$1/(2\sqrt2)$ is approximately $0.35355339$.
\end{theorem}

\subsection{The matrix and the numerical-radius formula}

On $\ell_4^3(\mathbb R)$ define
\[
 T=\frac1{1000}
 \begin{pmatrix}
 -10&2255&-40\\
 -1437&0&-1914\\
 34&2480&6
 \end{pmatrix}
 =\frac{A}{1000}.
 \tag{935.1}\label{eq:935-matrix}
\]
We will certify
\[
 v(T)\le\frac{5001}{5000},\qquad
 \|T\|_{4\to4}>\frac{3537}{1250}.
 \tag{935.2}\label{eq:935-bounds}
\]
For $x=(x_1,x_2,x_3)\in\mathbb R^3$, set
\[
 s(x)=x_1^4+x_2^4+x_3^4,
 \qquad F(x)=\sum_{i=1}^3x_i^3(Tx)_i.
\]
The unique norming functional of a unit vector of $\ell_4^3$ is
$(x_1^3,x_2^3,x_3^3)$, so
\[
 v(T)=\sup_{x\ne0}\frac{|F(x)|}{s(x)}.
 \tag{935.3}\label{eq:935-radius}
\]

\subsection{Two exact positive Gram matrices}

Let
\[
 z(x)=(x_1^2,x_2^2,x_3^2,x_1x_2,x_1x_3,x_2x_3)^{\mathsf T}
\]
and define the following symmetric integer matrices:
\[
 M_+=
 \begin{pmatrix}
10102&-6298&-14&-11275&200&214\\
-6298&10002&-7757&7185&-4&9570\\
-14&-7757&9942&-188&-170&-12400\\
-11275&7185&-188&12596&-214&4\\
200&-4&-170&-214&28&188\\
214&9570&-12400&4&188&15514
 \end{pmatrix},
 \tag{935.4}\label{eq:935-plus}
\]
\[
 M_-=
 \begin{pmatrix}
9902&-6433&-8&11275&-200&218\\
-6433&10002&-7646&-7185&-4&-9570\\
-8&-7646&10062&-204&170&12400\\
11275&-7185&-204&12866&-218&4\\
-200&-4&170&-218&16&204\\
218&-9570&12400&4&204&15292
 \end{pmatrix}.
 \tag{935.5}\label{eq:935-minus}
\]
Direct expansion gives the exact polynomial identities
\[
 \frac1{10000}z(x)^{\mathsf T}M_+z(x)
   =\frac{5001}{5000}s(x)-F(x),
 \tag{935.6}\label{eq:935-sosplus}
\]
\[
 \frac1{10000}z(x)^{\mathsf T}M_-z(x)
   =\frac{5001}{5000}s(x)+F(x).
 \tag{935.7}\label{eq:935-sosminus}
\]
For example, the coefficient of $x_1^3x_2$ in the numerator on the left
of \eqref{eq:935-sosplus} is $2(-11275)=-22550=-10A_{12}$.
All coefficients of absent monomials cancel.

Both matrices are positive definite. Their leading principal minors, in
orders one through six, are the following positive integers:
\[
 \begin{aligned}
 M_+:\;&10102,\quad 61375400,\quad 976429002,\\
       &6319567977,\quad 39156751988,\quad156153016840;
 \end{aligned}
 \tag{935.8}\label{eq:935-minorsplus}
\]
\[
 \begin{aligned}
 M_-:\;&9902,\quad57656315,\quad466262882,\\
       &4413501692,\quad8528707160,\quad32937577888.
 \end{aligned}
 \tag{935.9}\label{eq:935-minorsminus}
\]
Sylvester's criterion and \eqref{eq:935-sosplus}--\eqref{eq:935-sosminus}
therefore imply $|F(x)|\le(5001/5000)s(x)$ for every real $x$.
Equation~\eqref{eq:935-radius} proves the numerical-radius bound in
\eqref{eq:935-bounds}.

\subsection{An exact norm witness and comparison with \texorpdfstring{$M_4$}{M4}}

For $w=(185,1000,-179)^{\mathsf T}$ one has
\[
 Aw=(2260310,76761,2485216)^{\mathsf T},
\]
\[
 \sum_iw_i^4=1002197976306,\qquad
 \sum_i(Aw)_i^4=64248592946651712875988977.
\]
The integer comparison
\[
 \begin{split}
 625\sum_i(Aw)_i^4-14148^4\sum_iw_i^4
 &=916493558274854032306529\\
 &>0
 \end{split}
 \tag{935.10}\label{eq:935-witness}
\]
is exactly
$\|Tw\|_4^4>(3537/1250)^4\|w\|_4^4$, because
$14148/5=1000(3537/1250)$. This proves the second bound in
\eqref{eq:935-bounds} and hence
\[
 \frac{v(T)}{\|T\|_{4\to4}}
 <\frac{5001/5000}{3537/1250}
 =\frac{1667}{4716}.
 \tag{935.11}\label{eq:935-ratio}
\]

For $0\le t\le1$ the identity
\[
 1+t^4-2\sqrt2(t-t^3)=(t^2+\sqrt2t-1)^2\ge0
\]
has equality at $t=(\sqrt6-\sqrt2)/2$. Therefore
\[
 M_4=\max_{0\le t\le1}\frac{t-t^3}{1+t^4}=\frac1{2\sqrt2}.
\]
Finally, $4716^2-8\cdot1667^2=9544>0$, so
$1667/4716<1/(2\sqrt2)$. Thus \eqref{eq:935-ratio} disproves the
corresponding three-dimensional formula.

\subsection{Pointwise amplification to the exact Lebesgue space}

Set $\mathcal X=L_4([0,1];\ell_4^3(\mathbb R))$ and define
$(\mathcal T f)(t)=T(f(t))$. Pointwise integration gives
$\|\mathcal T\|\le\|T\|_{4\to4}$, while constant vector-valued functions
give the reverse inequality. Thus the norms are equal. For a unit vector
$f\in\mathcal X$, its norming functional is pointwise $(f_1^3,f_2^3,f_3^3)$.
Integrating \eqref{eq:935-sosplus}--\eqref{eq:935-sosminus} gives
\[
 \left|\int_0^1\sum_i f_i(t)^3(\mathcal T f)_i(t)\,dt\right|
 \le\frac{5001}{5000}\int_0^1\sum_i|f_i(t)|^4\,dt.
\]
Consequently $v(\mathcal T)\le5001/5000$; indeed constant functions show
$v(\mathcal T)=v(T)$.

The explicit surjective isometry from $\mathcal X$ onto the scalar real
Lebesgue space is
\[
 (Uf)(s)=3^{1/4}f_i(3s-i+1),\qquad
 \frac{i-1}{3}\le s<\frac{i}{3},\quad i=1,2,3.
\]
Transport $\mathcal T$ through $U$ and normalize by its operator norm.
Equation~\eqref{eq:935-ratio} now proves the assertion of
Theorem~\ref{thm:935} for the precise space in the catalogue.
This is a pointwise amplification, not a conditional-expectation
compression; the latter would not automatically preserve the numerical
radius.

\subsection{The dual exponent \texorpdfstring{$4/3$}{4/3}}

For an operator $R$ on a reflexive real Banach space $X$,
$\|R^*\|=\|R\|$ and $v(R^*)=v(R)$. To verify the latter identity, every
norming pair $(x,x^*)$ for $X$ gives the norming pair $(x^*,J_Xx)$ for
$X^*$, with
\[
 (J_Xx)(R^*x^*)=x^*(Rx).
\]
Reflexivity supplies the converse correspondence for every norming pair in
$X^*\times X^{**}$.
Apply this to $X=\ell_4^3$. The matrix $T^{\mathsf T}$ on
$\ell_{4/3}^3$ has the same norm and numerical radius as $T$.
Pointwise amplification works for every $1<p<\infty$: the norming functional
is $\operatorname{sign}(f_i)|f_i|^{p-1}$, and the finite-dimensional
numerical-radius inequality integrates with degree $p$.
Hence \eqref{eq:935-ratio} also applies on real $L_{4/3}([0,1])$.
The substitution $s=t^3$ in the defining maximum gives $M_{4/3}=M_4$.

\paragraph{Verification and scope.}
The displayed integers form a complete certificate. The supplied
\texttt{verify935\_exact.py} checks both polynomial identities, all twelve
principal minors, \eqref{eq:935-witness}, and the strict constant comparison
using only exact integer and rational arithmetic. An independent verifier
uses separate formulas for the fifteen quartic coefficients and rational
$LDL^{\mathsf T}$ decomposition. Floating-point searches were used only to
discover the certificate, and are not premises of the proof.
The two-dimensional theorem of \cite{Chiclana2026} is not contradicted:
this matrix has dimension three. The counterexample refutes the real-case conjecture in
\cite[Problem~5.1]{MMM2011} and resolves the exact catalogue identity
question \cite{AIM}; it does not determine the numerical index itself.

\subsection{A certified interval of exponents}
\label{sec:935-robustness}

The same matrix yields a quantitative neighborhood of counterexamples.
\begin{proposition}\label{prop:935-interval}
The identity in Problem~935 fails for every exponent in either interval
\[
 \left[\frac{39999}{10000},\frac{40001}{10000}\right],
 \qquad
 \left[\frac{40001}{30001},\frac{39999}{29999}\right].
\]
\end{proposition}
\begin{proof}
Write $v_p$ and $\|\cdot\|_p$ for the numerical radius and operator norm on
real $\ell_p^3$, and keep the matrix $T$ from \eqref{eq:935-matrix}. Set
\[
 \varepsilon=\frac1{10000},\quad
 \gamma=\frac{5001}{5000},\quad L=\frac{3537}{1250},\quad
 R=\|T\|_{\infty\to\infty}=\frac{3351}{1000}.
\]
We first treat $|p-4|\le\varepsilon$. For a nonzero real vector $x$, use
homogeneity to arrange $\max_i|x_i|=1$, and put
\[
 S_p(x)=\sum_i|x_i|^p,\qquad
 F_p(x)=\sum_i\operatorname{sign}(x_i)|x_i|^{p-1}(Tx)_i.
\]
Then $S_p(x)\ge1$, and at least one coordinate has absolute value one.
For $a>0$,
\[
 \max_{0\le t\le1}t^a|\log t|=\frac1{ea},
\]
with the value at zero understood by continuity. The mean value theorem,
$p\ge39/10$, and $|(Tx)_i|\le R$ therefore imply
\[
 |F_p(x)-F_4(x)|
 \le\frac{2R}{e(29/10)}\varepsilon,
 \qquad
 |S_p(x)-S_4(x)|
 \le\frac{2}{e(39/10)}\varepsilon.
 \tag{935.12}\label{eq:935-pvariation}
\]
The factor two is valid because the summand corresponding to a coordinate
of absolute value one vanishes in each difference. The certificate already
proved $|F_4(x)|\le\gamma S_4(x)$. Since $e>8/3$, it follows that
\[
 \begin{split}
 |F_p(x)|
 &\le\gamma S_p(x)+C\varepsilon,\\
 C&=\frac{15}{2}\left(\frac{R}{29}+\frac{\gamma}{39}\right)
   =\frac{399237}{377000}.
 \end{split}
\]
After division by $S_p(x)\ge1$ and taking the supremum,
\[
 v_p(T)\le\gamma+C\varepsilon.
 \tag{935.13}\label{eq:935-vp}
\]

Standard comparison of the three-dimensional $\ell_p$ norms gives
\[
 \|T\|_p\ge3^{-d}\|T\|_4,
 \qquad d=\left|\frac1p-\frac14\right|
          \le\frac{\varepsilon}{4p}.
\]
For $0\le d\le1$, convexity gives $3^d\le1+2d$. As $p\ge39/10>7/2$,
\[
 3^d\le1+\frac{\varepsilon}{2p}
       \le1+\frac{\varepsilon}{7}.
\]
The strict bound $\|T\|_4>L$ hence yields
\[
 \frac{v_p(T)}{\|T\|_p}
 <\frac{(\gamma+C\varepsilon)(1+\varepsilon/7)}{L}
 =\frac{87994832581079}{248910480000000}.
 \tag{935.14}\label{eq:935-rp}
\]

For $p>2$, put $f_p(t)=(t-t^{p-1})/(1+t^p)$ on $[0,1]$.
Direct differentiation gives
\[
 \frac{\partial f_p(t)}{\partial p}
 =\frac{(-\log t)t^{p-1}(1+t^2)}{(1+t^p)^2},
 \qquad
 \left|\frac{\partial f_p(t)}{\partial p}\right|
 \le\frac{2}{e(p-1)}\le\frac{15}{58}
\]
throughout the interval under consideration. Evaluate $f_p$ at the maximizing
point $t_0=(\sqrt6-\sqrt2)/2$ of $f_4$. Since
$8\cdot7071^2=399992328<400000000=20000^2$,
\[
 M_p\ge f_p(t_0)
 \ge M_4-\frac{15}{58}\varepsilon
 >\frac{7071}{20000}-\frac{15}{58}\varepsilon
 =\frac{51261}{145000}.
 \tag{935.15}\label{eq:935-Mp}
\]
The two rational bounds are strictly separated:
\[
 \frac{51261}{145000}
 -\frac{87994832581079}{248910480000000}
 =\frac{1030282921}{248910480000000}>0.
\]
Thus $v_p(T)/\|T\|_p<M_p$ throughout the first interval.
Pointwise amplification gives the same strict inequality for the numerical
index of real $L_p([0,1])$.

For the conjugate exponent $q=p/(p-1)$, reflexive duality gives
$v_q(T^{\mathsf T})=v_p(T)$ and
$\|T^{\mathsf T}\|_q=\|T\|_p$. The substitution $s=t^{p-1}$ in the
maximum defining $M_q$ gives $M_q=M_p$. The decreasing map
$p\mapsto p/(p-1)$ carries the first interval onto the second, proving
that case as well.
\end{proof}



\begin{thebibliography}{99}
\small
\setlength{\itemsep}{3pt}

\bibitem{Chiclana2026}
R.~Chiclana,
\emph{The Numerical Index of Two-Dimensional Real $\ell_p$ Spaces},
preprint, June 2026.
\href{https://arxiv.org/abs/2606.01675}{arXiv:2606.01675}.

\bibitem{MMM2011}
M.~Mart\'in, J.~Mer\'i, and M.~Popov,
\emph{On the numerical radius of operators in Lebesgue spaces},
J. Funct. Anal. \textbf{261} (2011), 149--168.
\href{https://doi.org/10.1016/j.jfa.2011.03.007}{doi:10.1016/j.jfa.2011.03.007}.
Problem 5.1.

\bibitem{AIM}
MColbrook/AIM, \emph{AIM 935: The numerical index of real Lebesgue spaces},
repository statement at revision \texttt{5d7e26e8444e}.
\url{https://github.com/MColbrook/AIM/blob/5d7e26e8444eb98de91988454564e45e011107c4/problems/935-the-numerical-index-of-real-lebesgue-spaces.md}.
\end{thebibliography}
\end{document}
